% ==================== КОДИРОВКА И ЯЗЫК ====================
\usepackage[T2A]{fontenc}
\usepackage[utf8]{inputenc}
\usepackage[russian]{babel}

% ==================== МНОГОФАЙЛОВЫЕ ПРОЕКТЫ ====================
\usepackage{subfiles}
\usepackage{standalone}

% ==================== ШРИФТ (Times-подобный, как в эталоне) ====================
\usepackage{newtxtext}
\usepackage{newtxmath}

% ==================== ГЕОМЕТРИЯ ====================
\usepackage{geometry}
\geometry{
  left=3cm,
  right=1.5cm,
  top=2cm,
  bottom=2cm,
  footskip=1cm
}

% ==================== ГРАФИКА И ДИАГРАММЫ ====================
\usepackage{tikz}
\usetikzlibrary{shapes.geometric, positioning, arrows.meta, calc,
                shadows, backgrounds, fit}
\usepackage{graphicx}
\usepackage{float}
\usepackage{caption}
\usepackage{pdflscape}

% ==================== МАТЕМАТИКА ====================
\usepackage{amsmath}

% ==================== ТЕКСТ ====================
\usepackage{setspace}
\onehalfspacing
\usepackage{indentfirst}
\usepackage{titlesec}
\usepackage{enumitem}
\usepackage{xurl}          % переносы длинных технических имён

% ==================== ТАБЛИЦЫ ====================
\usepackage{array}
\usepackage{longtable}
\usepackage{booktabs}

% ==================== ГИПЕРССЫЛКИ ====================
\usepackage[hidelinks]{hyperref}

% ==================== ЗАГОЛОВКИ ====================
% Section: 14pt жирный, слева, как в эталоне
\titleformat{\section}
  {\normalfont\bfseries\large}
  {\thesection}
  {1em}
  {}
\titleformat{\subsection}
  {\normalfont\bfseries\normalsize}
  {\thesubsection}
  {1em}
  {}

\titlespacing*{\section}{0pt}{1.5\baselineskip}{1\baselineskip}
\titlespacing*{\subsection}{0pt}{1.2\baselineskip}{0.8\baselineskip}

% ==================== ПОДПИСИ РИСУНКОВ И ТАБЛИЦ ====================
\captionsetup{
  font=normalsize,
  labelfont=normalfont,
  labelsep=colon,             % "Рис. 1: ...", "Таблица 1: ..."
  justification=centering,
  singlelinecheck=false
}

% ==================== ПАРАМЕТРЫ СТРАНИЦЫ ====================
\setlength{\parindent}{1.25cm}
\setlength{\parskip}{0pt}
\setlength{\emergencystretch}{3em}   % защита от переполнения строк
\tolerance=2000
\hyphenpenalty=1000
\sloppy                              % глобально разрешить "прослабленные" строки

% ==================== КАСТОМНЫЕ КОМАНДЫ ====================
\newcommand{\field}[1]{\texttt{#1}}
\urlstyle{tt}                        % \path{...} и \url{...} моноширинным

% Вставка скриншота: если файла ещё нет, выводится рамка-заглушка.
% \screenshot{путь}{подпись}{метка}{текст-подсказка}
\newcommand{\screenshot}[4]{%
  \begin{figure}[H]
    \centering
    \IfFileExists{#1}{%
      \includegraphics[width=\linewidth]{#1}%
    }{%
      \fbox{\begin{minipage}[c][6cm][c]{0.92\linewidth}
        \centering\small
        \textbf{Место для скриншота}\\[4pt]
        #4
      \end{minipage}}%
    }%
    \caption{#2}
    \label{#3}
  \end{figure}
}
